\chapter{Discussion}
\label{chap:discussion}

\section{What actually reduced gas}

The experimental results show that the main gas reduction did not come from a single
micro-optimization. It came from separating concerns that had been merged in the first
implementation of the variants. The intermediate monolithic contracts were useful because they
made all variants executable, but they also deployed mutually exclusive logic together. The
final implementation reduced gas by moving from this ``one contract with many flags'' model
to specialized contracts whose bytecode matches the selected protocol mode.

The clearest example is the optimistic path. In the initial implementation, each exchange
deployed a full optimistic account. In the final clone-based architecture, the expensive
implementation contracts are deployed once and each new exchange only creates a minimal
clone. This changes the gas model from a repeated full deployment to a fixed infrastructure
cost plus a low marginal cost. As shown in Chapter~\ref{chap:experimental-results}, the
measured marginal optimistic costs are \(500,720\) gas for \texttt{no\_S\_deposit},
\(516,349\) gas for \(S=B\), and \(556,309\) gas for \(S=V\). These are the strongest
honest-path gains of the project.

The second real reduction comes from removing protocol steps that are unnecessary in some
role configurations. The \(S=B\) case fuses the sponsor deployment/payment role with the
buyer role in the optimistic phase. The \(S_B=B, S_V=V\) dispute case removes the need to
repeat the challenge loop through Step 9 with another sponsor configuration. This is not a
low-level Solidity trick. It is a protocol-level simplification reflected in the
implementation. The measured variable dispute cost after the trigger is almost halved in
the self-sponsored case, which matches the expected reduction from avoiding repeated
Step 7--8 iterations.

The third reduction comes from the hardcoded \(\mathsf{SHA256}\) dispute path. Here, the
contract does not accept a generic circuit gate and a proof \(\pi_1\) as the source of
truth. Instead, it reconstructs the expected gate from the public metadata of the
hardcoded \(desc=\mathsf{SHA256}(x)\) circuit. This removes the generic
\(h_{\mathit{circuit}}\) membership proof from Step 8. Locally, this saves \(62,542\) gas
on the 16 KiB Step 8 measurement and about \(166,000\)--\(169,000\) gas on the isolated
1 GiB \(h_{\mathit{circuit}}\) check. In the final specialized implementation, the
hardcoded path also reduces the dispute deployer cost, which is why the full hardcoded
16 KiB dispute improves much more than the isolated \(\pi_1\) saving alone.

The main lesson is that the successful optimizations are structural. They either avoid
deploying unused bytecode, avoid executing unnecessary protocol loops, or remove a generic
proof obligation when the circuit is already determined by the precontract.

\section{What did not reduce gas as expected}

Several ideas were useful conceptually but did not immediately reduce gas in the first
implementation. The most important example is the intermediate monolithic contract. Adding
\texttt{no\_S\_deposit}, \(S=B\), \(S=V\), self-sponsored disputes, and hardcoded
\(\mathsf{SHA256}\) as flags inside the existing contracts made the behavior configurable,
but it increased deployment gas. The honest optimistic path rose from \(2,300,201\) gas in
the initial implementation to \(3,046,782\) gas in the intermediate monolithic implementation.
This result was not a measurement error: it was a consequence of deploying more code per
exchange.

The first hardcoded measurements also looked weaker than expected when they were read as
full-dispute totals. This was because the hardcoded change affects a specific part of
Step 8, while a complete dispute also includes dispute deployment, challenge rounds,
Merkle proofs for \(h_{\mathit{ct}}\) and \(hpre\), gate evaluation, and finalization. If
one measures only the full total, the removal of \(\pi_1\) is diluted by the remaining
costs. The local hardcoded saving exists, but it is not the only component of the full
dispute. This explains why the hardcoded optimization must be reported both locally and
end-to-end.

Another important point is that the generic Step 8 verifier was not fundamentally
optimized in the normal non-hardcoded path. In the 4 MiB same-scope comparison, the final
specialized architecture reduces the overall dispute initialization cost because the
deployment architecture is better, but the generic \texttt{submitCommitment} cost remains
close to the monolithic implementation. Therefore, the final implementation improves the
architecture around the verifier more than it improves the generic verifier itself.

Finally, ERC-4337 support should not be described as a gas optimization in this project.
The inherited account-abstraction path is useful to demonstrate sponsored interactions and
UserOperations, but it introduces operational components such as the bundler and
EntryPoint. The specialized clone measurements deliberately use a simpler direct-call path.
This separation avoids claiming that account abstraction itself reduces the protocol gas.
In this implementation, the gas reductions come from specialization and clone deployment,
not from ERC-4337.

\section{Cost relocation versus true optimization}

One subtle point in the measurements is the difference between relocating a cost and
removing a cost. The clone architecture relocates part of the optimistic deployment cost
from every exchange to a one-time infrastructure deployment. This is a real practical gain
when multiple exchanges share the same infrastructure, but it is not the same as making the
first deployment free. The factory and implementation contracts cost \(6,584,001\) gas in
total. That cost must be paid once before the marginal per-exchange savings become
available.

This is still meaningful for a protocol implementation. In the initial architecture, the
same optimistic bytecode is redeployed for every exchange. In the final architecture, the
implementation bytecode is deployed once and reused. The marginal cost is the relevant
quantity for a deployed service, while the fixed infrastructure cost is relevant for the
first setup. The report therefore separates these two costs instead of adding the factory
deployment to every exchange or ignoring it completely.

The hardcoded \(\mathsf{SHA256}\) mode is different. It does not merely move a cost from
on-chain to an earlier deployment step. It removes the need to verify \(\pi_1\) for
\(h_{\mathit{circuit}}\) in the hardcoded case, because the contract can reconstruct the
expected circuit gate. However, it does not remove the need for witnesses altogether. The
contract still needs authenticated values for ciphertext blocks and intermediate
accumulators. These values are provided externally and verified against \(h_{\mathit{ct}}\)
or \(hpre\). Thus, hardcoding removes a specific generic proof obligation; it does not make
Step 8 free.

The off-chain large-file work is another example. Moving precontract generation and dispute
witness computation to the native Rust CLI is not an on-chain gas optimization. It is a
practical scalability improvement. It makes 1 GiB experiments possible without relying on a
browser or a JavaScript/WASM memory model. The on-chain contract still verifies the
submitted commitments and witnesses; the CLI only makes it feasible to generate the data
needed by the protocol.

These distinctions matter because a fair comparison must compare the same scope. A result
can be correct and still misleading if it mixes a fixed deployment cost with a marginal
per-exchange cost, or a local Step 8 saving with a complete dispute total. The final
experimental chapter was structured around this lesson.

\section{Trade-offs of specialization}

Specialization reduces gas, but it also introduces engineering trade-offs. The final
architecture contains more contracts than the initial implementation: factory, clone
implementations, direct variants, normal dispute accounts, self-sponsored dispute accounts,
hardcoded dispute accounts, deployer libraries, and verifier libraries. This makes the
bytecode paid by each exchange smaller, but it makes the source tree larger and increases
the number of paths that must be tested.

This trade-off is acceptable in this project because the variants are protocol modes, not
arbitrary runtime options. The mode \texttt{no\_S\_deposit}, \(S=B\), and \(S=V\) is known
at precontract time. The hardcoded \(\mathsf{SHA256}\) mode is also fixed before dispute.
Therefore, there is no reason to deploy a contract that contains all possible branches for
one exchange. Specialized contracts match the protocol structure more closely than a single
universal account.

The trade-off is more delicate for modes that are only known after Step 4, such as
\(S_B=B\) and \(S_V=V\). These cannot be chosen at the initial precontract deployment in
the general case. The implementation therefore specializes them at dispute deployment time.
This is why the optimistic contract keeps enough logic to select the correct dispute
deployer when the dispute sponsor roles are known. This design follows the protocol timing:
precontract-time facts are specialized early, dispute-time facts are specialized later.

There is also a compatibility trade-off. The normal direct path preserves the inherited
ERC-4337-compatible behavior. The clone variants are optimized for direct-call gas
measurements and simplified execution. This means the final implementation does not claim
that every optimized variant has the same account-abstraction integration as the normal
path. This is a deliberate boundary: preserving all ERC-4337 logic inside every optimized
clone would reintroduce the monolithic overhead that the project was trying to remove.

The strongest argument in favor of specialization is that it also helped with mainnet code
size constraints. The original generic dispute contract and deployer were close to or above
the EVM bytecode size limit in some configurations. Splitting the hardcoded dispute account,
deployer, and circuit library produced components below the limit. This is an architectural
benefit as well as a gas benefit.

\section{Practicality of large-file fair exchange}

The project made large-file experiments concrete. The native 1 GiB benchmark executed a
complete hardcoded dispute path up to the terminal state. This is important because the
previous browser-oriented workflow was not a reliable way to handle such files. The final
benchmark shows that the protocol implementation can generate large precontract artifacts,
compute dispute witnesses, execute 26 challenge rounds, perform Step 8, and finalize.

The result should nevertheless be interpreted carefully. The 1 GiB measurement is a
controlled benchmark, not a production user experience. The file is processed by native
Rust tooling, the blockchain is a local Hardhat network, and the measured transaction time
does not include public-network latency or mempool behavior. The benchmark proves
feasibility of the implemented cryptographic and on-chain path; it does not prove that the
current UI is ready for production-scale 1 GiB disputes.

The off-chain times are also significant. Preparing the 1 GiB precontract took
22.42--48.70 seconds across repeated runs on the benchmark machine. Generating all
\(hpre_i\) responses for the forced-left dispute took 9.47--11.92 seconds. These values are acceptable for an
experimental proof of concept, but they are visible to users. They also show why the native
CLI is necessary: the computation is not just a small frontend task.

On-chain, the 1 GiB self-sponsored hardcoded path costs \(13,990,092\) gas in the measured
complete scenario. This is not negligible, but it remains a concrete executed path rather
than a theoretical estimate. With external sponsors and three Step 7--8 iterations, the
measured total including the optimistic path is \(29,532,277\) gas. This larger value shows
why self-sponsoring and avoiding unnecessary iterations matter in practice.

The practical conclusion is therefore balanced. The implementation now handles 1 GiB
experiments in a meaningful end-to-end way, but the current best path relies on the
hardcoded \(\mathsf{SHA256}\) specialization and native tooling. A fully generic,
browser-first, production-ready large-file fair exchange remains future work.

\section{Lessons learned from the implementation}

The first lesson is that implementation-level gas intuition is unreliable without
same-scope measurements. Before the measurements were reorganized, several numbers looked
contradictory because fixed infrastructure costs, marginal exchange costs, full dispute
totals, and local Step 8 costs were mixed together. Once the scopes were separated, the
results became coherent: monolithic flags increase deployment cost, clones reduce marginal
honest-path cost, self-sponsoring reduces repeated loops, and hardcoding removes a specific
generic proof.

The second lesson is that protocol variants should be aligned with when information becomes
available. The modes \(S=B\), \(S=V\), and \texttt{no\_S\_deposit} are determined at
precontract time, so they can be deployed as distinct optimistic contracts. The modes
\(S_B=B\) and \(S_V=V\) are known only after the buyer expresses dissatisfaction and the
dispute sponsors are selected, so they belong in the dispute deployment decision. This
timing distinction is important both for correctness and for gas optimization.

The third lesson is that security review can change the interpretation of performance
work. The hardcoded \(\mathsf{SHA256}\) path is not only a gas optimization; it also
reduces reliance on the generic circuit encoding for that particular description. However,
the review also identified inherited issues in the generic V2 circuit format and Merkle
accumulator. These issues do not invalidate the measured gas improvements, but they show
that a clean next version should update the commitment format, not only the gas
architecture.

The fourth lesson is that an implementation can be improved without hiding its limitations.
The final version is better than the initial implementation on the main same-scope
measurements, and it supports large-file benchmarks that were not practical in the browser.
At the same time, it is still a research prototype with local accounts, a local backend,
development-chain measurements, and an inherited generic encoding that deserves a new
version. Stating these boundaries explicitly makes the results more credible, not weaker.

Overall, the project moved the implementation from a functional but expensive
one-contract-per-exchange prototype toward a more realistic architecture: reusable
optimistic infrastructure, specialized modes, hardcoded dispute verification for
\(\mathsf{SHA256}\), and native tooling for large files. The remaining work is now clearer:
reduce the dominant dispute costs, clean the cryptographic encoding, and decide how far the
factory/registry architecture should go in a future version.
